% Short report. Hard limit: fewer than 4 A4 pages, so the target is 3.
\documentclass[10pt,a4paper]{article}
\usepackage[margin=2cm]{geometry}
\usepackage[utf8]{inputenc}
\usepackage[T5,T1]{fontenc}
\usepackage{booktabs,graphicx,tabularx,hyperref}
\title{Evaluating Lightweight Vision-Language and OCR Models\\on Occluded Document Images}
\author{TODO}
\date{}

\begin{document}
\maketitle

\section{Introduction}
Real documents carry ink marks, stamps, dust and correction fluid that hide part of the text.
We ask how well \emph{lightweight} text recognisers read such pages, and at what cost:
(i)~which model transcribes occluded pages most accurately,
(ii)~which occlusions hurt most, and which model degrades least, and
(iii)~how accuracy trades off against model size and inference time.
We compare three OCR pipelines and two vision-language models (VLMs), ordered from reading only what is
visible to predicting text from a language prior. All run zero-shot with their pretrained weights:
nothing is trained or fine-tuned, and every setting comes from the model's documentation, fixed before
the full run. Each model reads the same 3{,}000 test pages of OPRB~\cite{oprb}. We score every transcript
under three protocols that separate reading errors from output format and reading order, and compare
models page by page with bootstrap confidence intervals.

\section{Dataset Statistics}
OPRB~\cite{oprb} draws synthetic occlusions on English arXiv pages from DocBank and rotates every page by up
to $\pm5^\circ$. Its test set has 6{,}000 pages occluded on text regions (\texttt{text}) and 3{,}000 on
non-text regions (\texttt{nontext}). Models see only the image; the ground truth (GT), one token per line
with its box, colour, font and layout class, is read only for scoring. Five models on 9{,}000 pages exceed
our GPU budget, so we draw 3{,}000 pages at random (seed 42), stratified by region, occlusion type and
density in the original proportions (Table~\ref{tab:data}); a sample mean thus estimates the test-set mean
without weights. All models read the same pages.

\begin{table}[h]
\centering\small
\caption{Pages per occlusion type in the test set and in our sample. Black Ink, White Whitener,
White Burnt and Through Dust come in three densities (0.5\,\%, 1\,\%, 2\,\% of the page, in roughly
equal shares); Sim and Through Stamp have no density level.}
\label{tab:data}
\begin{tabular}{@{}llrrrrrrr@{}}
\toprule
Region & Pages & \shortstack[r]{Black\\Ink} & \shortstack[r]{White\\Whitener} & \shortstack[r]{White\\Burnt}
  & \shortstack[r]{Through\\Dust} & \shortstack[r]{Through\\Stamp} & Sim & Total \\
\midrule
\texttt{text}    & test set & 1{,}000 & 1{,}000 & 1{,}000 & 1{,}000 & 1{,}000 & 1{,}000 & 6{,}000 \\
                 & sample   & 334 & 333 & 333 & 334 & 333 & 333 & 2{,}000 \\
\texttt{nontext} & test set & 537 & 537 & 537 & 537 & 429 & 423 & 3{,}000 \\
                 & sample   & 178 & 179 & 180 & 178 & 143 & 142 & 1{,}000 \\
\bottomrule
\end{tabular}
\end{table}

\paragraph{What the ground truth contains.} Several properties of the GT shape how scores should be read.
First, DocBank marks lines and figures with placeholder tokens (\texttt{\#\#LTLine\#\#},
\texttt{\#\#LTFigure\#\#}) that no model can read from the image. In the test set they are 1.5\,\% of GT
tokens on \texttt{text} pages but 45\,\% on \texttt{nontext} pages, where one page holds 137{,}800. Without
them, a sampled page has a median of 516 words (\texttt{text}) and 455 words (\texttt{nontext}).
Second, the GT lists tokens along a whole printed line, so on a two-column page it alternates between
the columns. Two-column pages make up 17\,\% of \texttt{text} but 44\,\% of \texttt{nontext} pages in
the sample, so the two regions differ in layout as well as in where the occlusion falls. Third,
according to the dataset card the GT describes the clean page, so text hidden under an occlusion still
counts and the task is full-page recovery. A few pages have almost no text after removing placeholders
(3 of 3{,}000 under 10 characters), which makes their per-page CER unbounded; we therefore also report
median and micro-averaged CER. The test set contains six occlusion types; it has no \textit{Black Scribble}
or \textit{Mixed} pages, so scribbles are not covered.

\section{Model Descriptions}
Table~\ref{tab:models} lists the five models. We chose them so that each step from top to bottom relies
less on what is visible and more on a language prior, which lets the results explain \emph{why} a model
is or is not robust.

\begin{table}[h]
\centering\small
\caption{Evaluated models (released weights, pinned versions). Parameters and weight sizes as loaded.
VLMs decode greedily with at most 4{,}096 new tokens; a page that repeats one pattern for its last 256
tokens is stopped and marked \emph{truncated}.}
\label{tab:models}
\begin{tabularx}{\textwidth}{@{}llrrX@{}}
\toprule
Model & Recognition & Params & Weights & Hardware and settings \\
\midrule
Tesseract 5.5 (\texttt{tessdata\_best}) & LSTM per line & -- & 15\,MB & CPU; $3\times$ upscale, deskew; automatic layout \\
PP-OCRv5 mobile (English) & detector + CTC per line & -- & 13\,MB & CPU (ONNX); shared reading order \\
docTR FAST-base + PARSeq & detector + AR decoder per word & 34\,M & 138\,MB & T4, fp32; rotated boxes; shared reading order \\
LightOnOCR-2-1B & VLM for OCR & 1.0\,B & 2.8\,GB & T4; fp16 LM, fp32 vision encoder; long edge 1{,}540\,px; no prompt \\
Qwen3-VL-2B-Instruct & general VLM & 2.1\,B & 4.3\,GB & T4, fp16; at most 2{,}048 visual tokens; transcription prompt \\
\bottomrule
\end{tabularx}
\end{table}

\textbf{Tesseract~5} is the classical CPU baseline. Its guide asks for at least 300\,DPI and deskewed pages;
OPRB pages are about 95\,DPI, so we upscale by 3 and deskew. \textbf{PP-OCRv5 mobile} is a very small deep
pipeline whose CTC recogniser reads each detected line from its pixels alone. \textbf{docTR} pairs the
FAST-base detector with PARSeq, whose autoregressive decoder adds context within a word. For both detector
pipelines one shared routine deskews the boxes, groups them into lines, splits columns and orders them.
\textbf{LightOnOCR-2-1B} is trained end to end for page OCR and writes Markdown, with LaTeX for formulas; its
card asks for a 1{,}540\,px longest edge. In fp16 its vision encoder produced NaN on a T4, so the encoder runs
in fp32. \textbf{Qwen3-VL-2B-Instruct} is a general VLM with the strongest language model of the five; the
prompt asks for all visible text in reading order, without describing the page or inventing unreadable text.

\section{Experimental Results}
% 1.6 pages.
\subsection{Setup}
% Hardware, scoring protocols, number of pages, failures and truncations, 95% bootstrap intervals.
% Released weights, zero-shot, settings fixed in advance, none chosen by score.
\subsection{Quantitative results}
% Table 1 shows both protocols side by side: `notebook` (all GT tokens) and `normalized`
% (without ##LT...## placeholders and output markup). The analysis uses `normalized`.
\subsection{Text-region vs. non-text-region occlusions}
\subsection{Robustness to occlusion types}
\subsection{Qualitative examples}
\subsection{Efficiency}
\subsection{Limitations}

\section{Conclusion}
% 0.25 page: which model for which situation.

\begin{thebibliography}{1}\small
\bibitem{oprb} OPRB: Occluded Pages Restoration Benchmark. Dataset:
\url{https://huggingface.co/datasets/kpurkayastha/OPRB}; paper: arXiv:2604.10077.
\end{thebibliography}

\end{document}
